\chapter{Conclusiones}

El desarrollo del proyecto permitió implementar exitosamente un mismo algoritmo utilizando distintos paradigmas de computación concurrente y paralela.

La implementación secuencial proporcionó una línea base confiable para validar la exactitud funcional del algoritmo.

Las estrategias multicore y OpenMP demostraron ser alternativas eficientes para arquitecturas de memoria compartida, ofreciendo reducciones significativas en el tiempo de ejecución con un esfuerzo moderado de implementación.

La implementación MPI evidenció las ventajas del paradigma distribuido para escenarios de mayor escala, aunque introdujo costos asociados a la comunicación entre procesos.

Por su parte, la implementación basada en CUDA puso de manifiesto el potencial del paralelismo masivo proporcionado por las GPU, especialmente para cargas de trabajo altamente paralelizables.

En términos generales, los resultados obtenidos permiten concluir que la selección de una estrategia de paralelización debe considerar simultáneamente:

\begin{itemize}
    \item Las características del problema.
    \item La infraestructura disponible.
    \item El costo de implementación.
    \item El desempeño esperado.
\end{itemize}

Finalmente, el proyecto permitió fortalecer competencias relacionadas con el diseño, implementación y evaluación de soluciones concurrentes y distribuidas, evidenciando la importancia de estas técnicas en el desarrollo de aplicaciones de alto desempeño.
